% GENERATED FILE. Edit canonical lessons, then run npm run generate:books.
% canonical-source-hash: fnv1a64:98d7aa3039ae8f2a
% canonical-lessons: GE-C174-abholen, GE-C174-abfahren, GE-C174-ankommen, GE-C174-ausziehen, GE-C174-aufraeumen

\chapter{To pick up, To depart, To arrive, To take off, To tidy up}
\label{ch:ge-to-pick-up-to-depart-to-arrive-to-take-off-to-tidy-up}
\hlchaptermodality{174}

\begin{chapteropening}
\textbf{By the end of this chapter:} \emph{I can say to pick up, to depart, to arrive, to take off and to tidy up.}

Complete the last lesson of chapter 174: I can say to pick up, to depart, to arrive, to take off and to tidy up.
\end{chapteropening}

\section[abholen]{abholen — to pick up, to collect}

\label{lesson:GE-C174-abholen}



\begin{quote}
Before the new one: say the German for together, then the German for alone.
\end{quote}

\subsection*{You'll want to know: abholen}
\textbf{abholen} — \textquotedblleft{}to pick up, to collect\textquotedblright{}.

\begin{cousinweb}
\textbf{abholen} is \emph{ab-}, \textquotedblleft{}off, away\textquotedblright{}, + \emph{holen}, \textquotedblleft{}to fetch\textquotedblright{} $\leftarrow$ Middle High German \emph{holen} (EtymDB 2.1): to collect and take away.
\end{cousinweb}

\begin{grammarlens}[title={What you've built}]
One more word for this chapter.
\end{grammarlens}

\subsection*{Your turn}
\begin{itemize}
\raggedright
  \item \emph{Say it:} \emph{abholen}
  \item \emph{Say it:} \emph{abholen}, once more
  \item \emph{Say it:} say \emph{abholen}
  \item \emph{Recall:} say the German for together, then the German for alone, then say \emph{abholen} again
\end{itemize}

\begin{tcolorbox}[breakable,colback=teal!4,colframe=teal!35!black,title={Before you move on}]
What does abholen mean? (\textquotedblleft{}To pick up, to collect\textquotedblright{}.) Say it once more.
\end{tcolorbox}
\section[abfahren]{abfahren — to depart}

\label{lesson:GE-C174-abfahren}



\begin{quote}
Before the new one: say the German for alone, then the German for to pick up.
\end{quote}

\subsection*{You'll want to know: abfahren}
\textbf{abfahren} — \textquotedblleft{}to depart\textquotedblright{}.

\begin{cousinweb}
\textbf{abfahren} is \emph{ab-}, \textquotedblleft{}off, away\textquotedblright{}, + \emph{fahren}, \textquotedblleft{}to travel\textquotedblright{}, which EtymDB 2.1 lists beside English \emph{fare}.
\end{cousinweb}

\begin{grammarlens}[title={What you've built}]
One more word for this chapter.
\end{grammarlens}

\subsection*{Your turn}
\begin{itemize}
\raggedright
  \item \emph{Say it:} \emph{abfahren}
  \item \emph{Say it:} \emph{abfahren}, once more
  \item \emph{Say it:} say \emph{abfahren}
  \item \emph{Recall:} say the German for alone, then the German for to pick up, then say \emph{abfahren} again
\end{itemize}

\begin{tcolorbox}[breakable,colback=teal!4,colframe=teal!35!black,title={Before you move on}]
What does abfahren mean? (\textquotedblleft{}To depart\textquotedblright{}.) Say it once more.
\end{tcolorbox}
\section[ankommen]{ankommen — to arrive}

\label{lesson:GE-C174-ankommen}



\begin{quote}
Before the new one: say the German for to pick up, then the German for to depart.
\end{quote}

\subsection*{You'll want to know: ankommen}
\textbf{ankommen} — \textquotedblleft{}to arrive\textquotedblright{}.

\begin{cousinweb}
\textbf{ankommen} is \emph{an-}, \textquotedblleft{}at\textquotedblright{}, + \emph{kommen}, \textquotedblleft{}to come\textquotedblright{} $\leftarrow$ Old High German \emph{kweman} (EtymDB 2.1). EtymDB lists Dutch \emph{aankomen} beside it.
\end{cousinweb}

\begin{grammarlens}[title={What you've built}]
One more word for this chapter.
\end{grammarlens}

\subsection*{Your turn}
\begin{itemize}
\raggedright
  \item \emph{Say it:} \emph{ankommen}
  \item \emph{Say it:} \emph{ankommen}, once more
  \item \emph{Say it:} say \emph{ankommen}
  \item \emph{Recall:} say the German for to pick up, then the German for to depart, then say \emph{ankommen} again
\end{itemize}

\begin{tcolorbox}[breakable,colback=teal!4,colframe=teal!35!black,title={Before you move on}]
What does ankommen mean? (\textquotedblleft{}To arrive\textquotedblright{}.) Say it once more.
\end{tcolorbox}
\section[ausziehen]{ausziehen — to take off (clothes), to move out}

\label{lesson:GE-C174-ausziehen}



\begin{quote}
Before the new one: say the German for to depart, then the German for to arrive.
\end{quote}

\subsection*{You'll want to know: ausziehen}
\textbf{ausziehen} — \textquotedblleft{}to take off (clothes), to move out\textquotedblright{}.

\begin{cousinweb}
\textbf{ausziehen} is \emph{aus-}, \textquotedblleft{}out\textquotedblright{}, + \emph{ziehen}, \textquotedblleft{}to pull\textquotedblright{} $\leftarrow$ Proto-Germanic \emph{*teuhaną} (EtymDB 2.1): to pull off, or to pull out of a house.
\end{cousinweb}

\begin{grammarlens}[title={What you've built}]
One more word for this chapter.
\end{grammarlens}

\subsection*{Your turn}
\begin{itemize}
\raggedright
  \item \emph{Say it:} \emph{ausziehen}
  \item \emph{Say it:} \emph{ausziehen}, once more
  \item \emph{Say it:} say \emph{ausziehen}
  \item \emph{Recall:} say the German for to depart, then the German for to arrive, then say \emph{ausziehen} again
\end{itemize}

\begin{tcolorbox}[breakable,colback=teal!4,colframe=teal!35!black,title={Before you move on}]
What does ausziehen mean? (\textquotedblleft{}To take off (clothes), to move out\textquotedblright{}.) Say it once more.
\end{tcolorbox}
\section[aufräumen]{aufräumen — to tidy up}

\label{lesson:GE-C174-aufraeumen}



\begin{quote}
Before the new one: say the German for to arrive, then the German for to take off.
\end{quote}

\subsection*{You'll want to know: aufräumen}
\textbf{aufräumen} — \textquotedblleft{}to tidy up\textquotedblright{}.

\begin{cousinweb}
\textbf{aufräumen} is \emph{auf} + the verb \textquotedblleft{}to make room\textquotedblright{}, Old High German \emph{rūmen} $\leftarrow$ Proto-Germanic \emph{*rūmijaną}, built on \emph{*rūmaz}, \textquotedblleft{}roomy\textquotedblright{} (EtymDB 2.1). English \emph{room} shares the root.
\end{cousinweb}

\begin{grammarlens}[title={What you've built}]
One more word for this chapter.
\end{grammarlens}

\subsection*{Your turn}
\begin{itemize}
\raggedright
  \item \emph{Say it:} \emph{aufräumen}
  \item \emph{Say it:} \emph{aufräumen}, once more
  \item \emph{Say it:} say \emph{aufräumen}
  \item \emph{Recall:} say the German for to arrive, then the German for to take off, then say \emph{aufräumen} again
\end{itemize}

\begin{tcolorbox}[breakable,colback=teal!4,colframe=teal!35!black,title={Before you move on}]
What does aufräumen mean? (\textquotedblleft{}To tidy up\textquotedblright{}.) Say it once more.
\end{tcolorbox}
